% ============================================================================
% templates/report/nota-docente/nota-docente.tex — plantilla: nota de docente (guion de clase)
% ----------------------------------------------------------------------------
% Tu guía para explicar una clase: objetivos, guion paso a paso, puntos clave,
% ejemplos, errores comunes y cierre. SIN diseño: lo pone academic-report.
% Uso: scripts/build.sh nota-docente.tex   (LuaLaTeX)
% ============================================================================
\documentclass{academic-report}
\tipodocumento{Nota de docente}
\universidad{Universidad Nacional de San Cristóbal de Huamanga}
\escuela{Escuela Profesional de Economía}
\curso{Nombre del curso}\docente{Edison Achalma B.Sc. Econ.}\periodo{2026-I}
\titulodocumento{Sesión NN — Título de la sesión}

\begin{document}
\encabezado

\section{Objetivos de la sesión}
\begin{itemize}
  \item Que el estudiante comprenda\ldots
  \item Que el estudiante aplique\ldots
\end{itemize}

\section{Guion de la clase}
\begin{description}
  \item[00–10 min · Apertura.] Pregunta motivadora y conexión con la sesión previa.
  \item[10–35 min · Desarrollo.] Explicación del concepto central con la pizarra.
  \item[35–60 min · Aplicación.] Ejercicio guiado y trabajo en parejas.
  \item[60–70 min · Cierre.] Síntesis y anuncio de la tarea.
\end{description}

\section{Puntos clave para explicar}
\begin{definicion}[Concepto central]
Enunciado preciso del concepto y su intuición.
\end{definicion}
\begin{ejemplo}
Ejemplo aplicado al contexto peruano.
\end{ejemplo}

\section{Errores comunes}
\begin{importante}
Confusión típica del estudiante y cómo anticiparla.
\end{importante}

\section{Cierre y tarea}
Idea de cierre y consigna de la tarea para la próxima sesión.

\end{document}
